\documentclass[aspectratio=169]{beamer}
\usetheme{Madrid}
\usecolortheme{beaver}

\usepackage[UTF8]{ctex}
\usepackage{amsmath, amssymb}
\usepackage{booktabs}
\usepackage{ragged2e}

\title[0311report]{report}
\date{2026-03-11}

\begin{document}

\begin{frame}
    \titlepage
\end{frame}

\section{问题与目标}

\begin{frame}{1月版本的核心问题}
    \begin{itemize}
        \item \textbf{没做平行世界，影响归因分析}
        \item \textbf{冲击系数固定}：正or负面对个体的冲击幅度几乎保持为常数，不因不同个体/事件等而改变
        \item \textbf{问卷太简单}
        \item \textbf{机制闭环不完整}：状态变化对后续行为路径影响弱
        \item \textbf{触发偏概率注入}
    \end{itemize}
    \vspace{0.3cm}
\end{frame}

\begin{frame}{初步改造}
    \begin{itemize}
        \item 平行世界
        \item 状态更新改为动态变化（环境+个体特征+当前step共同决定）
        \item 问卷扩展为更多维(acceptance / trust / helplessness / behavior)
        \item 状态变化影响后续行为路径
        \item 触发改为多段判定+审计日志，降低无依据触发
    \end{itemize}
\end{frame}

\begin{frame}{初步改造示意图}
    \centering
    \includegraphics[width=\textwidth,height=0.82\textheight,keepaspectratio]{\detokenize{/Users/pifazuoren/Downloads/IMG_5454 2.jpg}}
\end{frame}

\section{机制升级}

\begin{frame}{触发机制完善}
    \begin{enumerate}
        \item 抽取当步信号：\texttt{step\_outcome / consumed\_time / failure\_pressure / success\_support}
        \item 数字语境检测：\texttt{检测action or status or intention }
        \item llm主导判定：\texttt{产出数字性+场景+置信度}
        \item 规则兜底：\texttt{first gate + status gate}
        \item 回退：\texttt{上述均未命中时才进入}
    \end{enumerate}
    \vspace{0.2cm}
    每次 attempt 都写入 \texttt{decision\_json}，存放路径来源、置信度、放行理由、\texttt{roll}、\texttt{event\_total\_prob}，便于审查。
\end{frame}



\begin{frame}{Hybrid校准与历史惯性}
    \begin{itemize}
        \item Hybrid 不让 LLM 直接决定是否数字，而是输出倍率：\texttt{risk\_mult/protect\_mult/neg\_impact\_mult/pos\_impact\_mult}
        \item 系统按 \(\alpha\) 融合 rule 与 LLM，低置信/跳变过大自动回退 rule
        \item attempt-hazard 引入短期记忆：近期事件密度、失败压力变化、refractory 衰减
    \end{itemize}
    \vspace{0.2cm}
    \[
    z_{total}=0.45 z_{prior}+0.55 z_{context}+offset-refractory+0.16\Delta fail+0.08\Delta mean
    \]
    \[
    p_{total}=\sigma(z_{total}),\quad
    p^- = p_{total}\cdot r_{neg},\quad
    p^+ = p_{total}(1-r_{neg})
    \]
\end{frame}

\begin{frame}{Prompt Engineering改造前后对比}
    \begin{center}
        \small
        \begin{tabular}{p{3.3cm}p{2.8cm}p{2.8cm}}
            \toprule
            指标 & 改造前 & 改造后 \\
            \midrule
            Hybrid 生效次数 & \(0\)（均回退为rule） & \(126\)（均为hybrid） \\
            非数字事件占比 & \(23/53=43.4\%\) & \(0/83=0\%\) \\
            高质量触发样本 & \(27/53=50.9\%\) & \(78/83=94\%\) \\
            \bottomrule
        \end{tabular}
    \end{center}
    \vspace{0.2cm}
    \begin{itemize}
        \item 样本纯度与可解释性显著提升。
    \end{itemize}
\end{frame}

\section{初步实验}

\begin{frame}{实验配置与规模}
    \begin{itemize}
        \item Group ID：\texttt{digital\_friction\_parallel\_20260311\_002952}
        \item 世界：A=\texttt{baseline\_low\_friction}，B=\texttt{high\_friction\_low\_assist}，C=\texttt{low\_friction\_high\_assist}
        \item 配对：4 组 seed（101--104），每组 A/B/C 同 seed
        \item 总子实验数：\(4\times 3=12\)
        \item 每次子实验：6个agents，单阶段1天，30分钟步长
        \item 运行总时长约 7.4 小时（平均每子实验约 37 分钟）
    \end{itemize}
\end{frame}


\section{结果}

\begin{frame}{行为指标结果（4组seed平均）}
    \begin{center}
        \small
        \begin{tabular}{lccc}
            \toprule
            World & AttemptRate & EmitGivenAttempt & NegShare \\
            \midrule
            A Baseline & 0.3637 & 0.2608 & 0.2640 \\
            B High Friction & 0.3429 & 0.5727 & 0.5402 \\
            C Low Friction + Assist & 0.2856 & 0.1505 & 0.1116 \\
            \bottomrule
        \end{tabular}
    \end{center}
    \vspace{0.2cm}

\end{frame}

\begin{frame}{心理状态变化（末值-初值，4组seed平均）}
    \begin{center}
        \small
        \begin{tabular}{lccc}
            \toprule
            World & Helplessness \(\Delta\) & Trust \(\Delta\) & Avoidance \(\Delta\) \\
            \midrule
            A Baseline & -5.58 & +1.73 & -2.19 \\
            B High Friction & +33.08 & -40.61 & +29.86 \\
            C Low Friction + Assist & -8.50 & +7.37 & -4.35 \\
            \bottomrule
        \end{tabular}
    \end{center}
    \vspace{0.2cm}

\end{frame}

\begin{frame}{Survey扩展指标（末轮-首轮，4组seed平均）}
    \begin{center}
        \scriptsize
        \begin{tabular}{lccc}
            \toprule
            指标 & A Baseline & B High Friction & C Low Friction + Assist \\
            \midrule
            TechAcceptance\(\Delta\) & +3.13 & +0.63 & -2.92 \\
            SurveyTrust\(\Delta\) & -0.83 & +0.63 & -4.79 \\
            SurveyAvoidance\(\Delta\) & +0.42 & -3.33 & -4.38 \\
            HelplessIndex\(\Delta\) & +3.13 & +0.31 & -1.56 \\
            WithdrawalIndex\(\Delta\) & -1.81 & +5.90 & -0.35 \\
            DelayOnline\(\Delta\) & -1.88 & +5.00 & -4.38 \\
            ProxyReliance\(\Delta\) & -6.04 & +8.33 & +1.67 \\
            OfflineSwitch\(\Delta\) & +2.50 & +4.38 & +1.67 \\
            \bottomrule
        \end{tabular}
    \end{center}
    \vspace{0.15cm}
    \tiny 注：本页由 \texttt{agent\_survey} 两轮问卷直接计算（末轮-首轮），用于补充展示。
\end{frame}

\begin{frame}{Survey扩展指标配对差值（B-A, C-A）}
    \begin{center}
        \scriptsize
        \begin{tabular}{lcc}
            \toprule
            指标 & B-A 均值差 & C-A 均值差 \\
            \midrule
            TechAcceptance\(\Delta\) & -2.50 & -6.04 \\
            SurveyTrust\(\Delta\) & +1.46 & -3.96 \\
            SurveyAvoidance\(\Delta\) & -3.75 & -4.79 \\
            HelplessIndex\(\Delta\) & -2.81 & -4.69 \\
            WithdrawalIndex\(\Delta\) & +7.71 & +1.46 \\
            DelayOnline\(\Delta\) & +6.88 & -2.50 \\
            ProxyReliance\(\Delta\) & +14.38 & +7.71 \\
            OfflineSwitch\(\Delta\) & +1.88 & -0.83 \\
            \bottomrule
        \end{tabular}
    \end{center}
\end{frame}



\begin{frame}{配对差值统计（B-A, C-A）}
    \begin{center}
        \small
        \begin{tabular}{lcc}
            \toprule
            指标 & B-A 均值差 & C-A 均值差 \\
            \midrule
            NegShare & +0.2762 & -0.1524 \\
            Helplessness\(\Delta\) & +38.66 & -2.92 \\
            Trust\(\Delta\) & -42.33 & +5.64 \\
            Avoidance\(\Delta\) & +32.06 & -2.15 \\
            \bottomrule
        \end{tabular}
    \end{center}
    \vspace{0.2cm}
\end{frame}

\section{结论与下一步}


\begin{frame}{下一步计划}
    \begin{itemize}
        \item 增加子实验次数
        \item 对于参数，做敏感性分析（步长、alpha、hazard窗口）
        \item 进一步校准初始化与冲击系数（这个感觉很难。。）
    \end{itemize}
\end{frame}

\begin{frame}
    \centering
    \Huge{请老师批评指正orz}
\end{frame}

\end{document}
